\documentclass[10pt,a4paper]{amsart}
\usepackage[margin=19mm]{geometry}
\usepackage{amsmath,amssymb,mathtools}
\usepackage[scr=rsfs,cal=euler]{mathalfa}
\usepackage{xfrac}
\usepackage[unicode,hidelinks,bookmarks=false]{hyperref}
\newcommand{\ie}{i.e.\ }
\newcommand{\cf}{cf.\ }
\newcommand{\resp}{respectively\ }
\newcommand{\Subsection}{\subsection}
\providecommand{\nobreakdash}{\nobreak\hbox{-}}
\setlength{\emergencystretch}{2em}
\multlinegap0pt
\usepackage{bm}
\usepackage{bbm}
\usepackage{dsfont}
\usepackage{xspace}
\usepackage{cancel}
\usepackage{mathtools}
\usepackage{amsthm}
\makeatletter
\@ifundefined{lemma}{\newtheorem{lemma}{Lemma}}{}
\@ifundefined{theorem}{\newtheorem{theorem}[lemma]{Theorem}}{}
\makeatother
\DeclareMathOperator{\HH}{H}
\DeclareMathOperator{\Hom}{Hom}
\DeclareMathOperator{\id}{id}
\newcommand{\Fp}{\mathbb{F}_p}
\newcommand{\F}{\mathbb{F}}
\newcommand{\Z}{\mathbb{Z}}
\newcommand{\g}{\mathfrak{g}}
\newcommand{\isom}{\cong}
\newcommand{\kk}{\mathfrak{k}}
\newcommand{\q}{\mathfrak{q}}
\title[Cohomology of solvable saturable pro-$p$ groups and Lie alge]{Cohomology of solvable saturable pro-$p$ groups and Lie algebras}
\author{Oihana Garaialde Oca\~na, Jon Gonz\'alez-S\'anchez, Lander Guerrero-S\'anchez}
\date{Source: arXiv:2507.18163v1 (CC BY 4.0)}
\begin{document}
\maketitle

\noindent\emph{Source and licence.} Cohomology of solvable saturable pro-$p$ groups and Lie algebras. Version arXiv:2507.18163v1 (CC BY 4.0), distributed under the Creative Commons Attribution 4.0 International licence, as recorded in the arXiv licence metadata for this exact version and shown on the abstract page https://arxiv.org/abs/2507.18163v1. Full rights notice: licenses/arxiv-2507.18163-CC-BY-4.0.txt.\par\smallskip

\noindent\textit{Editorial scope.} Counted unit: the Theorem whose source label is \texttt{None} in the article (source0.tex lines 536--539, source label \texttt{None}), delivered together with the article's own complete original proof of it (source0.tex lines 541--570, one \texttt{proof} environment of 30 lines). No mathematics is written, repaired or re-derived: statement and proof are the article's own text line for line. Required context retained uncounted: lem: abelian (lines 490--495). Every definition, display and label of those lines is retained unchanged. Omitted from this package and not redistributed: 1--489, 496--535, 540--540, 571--635. Nothing inside a retained excerpt is omitted or altered: every retained body is delivered exactly as the article's lines are written, so no omission receipt of any kind accompanies this package. Citations: the retained excerpts carry no citation command at all, so no bibliography entry of the article is transcribed and no reference is redistributed. Reference adaptation: none was needed and none was made; every reference inside the delivered unit resolves inside it, so no unresolved reference or citation is produced. Printed slips kept literal: no printed word, symbol, hypothesis or number of the counted statement or of its proof was altered. Layout and class: the article's own class is \texttt{article} with packages including a4wide, babel, amsmath, amsthm, amsfonts, amssymb, enumitem, mathtools, graphicx, wrapfig, tikz-cd, hyperref; the wrapper is the lane's pinned \texttt{amsart} editorial wrapper at 10pt on A4, which restates the theorem-like environments the retained text needs and adds the article's own macro definitions that the retained text needs (\texttt{\textbackslash F}, \texttt{\textbackslash Fp}, \texttt{\textbackslash HH}, \texttt{\textbackslash Hom}, \texttt{\textbackslash Z}, \texttt{\textbackslash g}, \texttt{\textbackslash id}, \texttt{\textbackslash isom}, \texttt{\textbackslash kk}, \texttt{\textbackslash q}), with extra wrapper packages (bm, bbm, dsfont, xspace, cancel, mathtools). The delivered numbering is the unit's own. Assets: no retained span contains a figure environment, a graphics-inclusion command, a PDF-page inclusion, a table or a TikZ picture; no image file is copied into this package, the package declares no asset dependency and no image file is redistributed with it. Licence: version arXiv:2507.18163v1 of this article is distributed under the Creative Commons Attribution 4.0 International licence, as recorded in the arXiv licence metadata for this exact version and shown on the abstract page https://arxiv.org/abs/2507.18163v1. A full rights notice is delivered at licenses/arxiv-2507.18163-CC-BY-4.0.txt. Characters: every retained excerpt is pure ASCII, so the delivered engine, XeLaTeX with its default Latin Modern text font, reports no missing character.
		\begin{lemma}\label{lem: abelian} Let $Q\cong \Z_p$ and let $\q\cong \Z_p$ be its associated $\Z_p$-Lie algebra. Let $V$ be an $\F_p$-vector space such that the action of $Q$ on $V$ is unipotent of class at most $p$. Then, for all $n\geq 0$, we have an $Q$-equivariant isomorphism 
		\begin{displaymath}
			\HH^n(Q;V) \isom \HH^n(\q; V)
		\end{displaymath}
		of  $\q$-modules, where the action of $\q$ on $\HH^n(\q; V)$ is given by the $\log$ functor.
	\end{lemma} 

	\begin{theorem}  
		Under the hypothesis described above, for all $r,s\geq 0$, we have that $E^{r,s}_2(G)\cong E_2^{r,s}(\g)$, and therefore, $E^{r,s}_{\infty}(G)\cong E^{r,s}_{\infty}(\g)$ as $G/K$-modules.

	\end{theorem}

	\begin{proof}
	
	Assume by induction that, for all $n\geq 0$, there is a $G$-equivariant isomorphism $$\HH^n(K;\Fp)\isom\HH^n(\kk;\Fp).$$
	Then, we have that 
	\begin{displaymath}
		E_2^{r,s}(G)=\HH^r\big(G/K;\HH^s(K,\Fp)\big)\isom \begin{dcases}
			\HH^s(K;\Fp)^{G/K} & \text{if } r=0, \\
			\frac{\HH^s(K;\Fp)}{(\sigma-1)\HH^s(K;\Fp)} & \text{if } r=1, \\
			0 & \text{if } r\geq 2, \\
		\end{dcases}
	\end{displaymath}
	and, similarly,
	\begin{displaymath}
		E_2^{r,s}(\g)=\HH^r\big(\g/\kk,\HH^s(\kk,\Fp)\big)\isom \begin{dcases}
			\HH^s(\kk,\Fp)^{\g/\kk} & \text{if } r=0, \\
			\frac{\HH^s(\kk,\Fp)}{T\HH^s(\kk,\Fp)} & \text{if } r=1, \\
			0 & \text{if } r\geq 2. \\
		\end{dcases}
	\end{displaymath}
To show that for all $r,s\geq 0$, Lazard's correspondence gives an isomorphism of $\F_p$-vector spaces $E^{r,s}_2(G)\cong E^{r,s}_2(\g)$, we set $V=\HH^s(K,\Fp)\isom \HH^s(\kk,\Fp)$ and verify that the conditions of Lemma \ref{lem: abelian} are satisfied. Indeed, $G$ acts on $\kk$ and thus, on $\kk/p\kk$, by the $\exp$ of the adjoint. Moreover, since $G$ is saturable, we have that $\gamma_p(G)\leq G^p$ and  thus, the action of $G$ on $\kk/p\kk$ is unipotent of class at most $p$. More explicitly, the generator $T=\log \sigma\in \q=\g/\kk$ acts on $x\in \kk$ via
	\begin{displaymath}
		x\cdot T=\sum_{i=1}^{p-1}\frac{(-1)^{i+1}}{i}x(C_\sigma-\id_Q)^i =x\cdot\log(C_{\sigma}),
	\end{displaymath}
	where $xC_\sigma=x^\sigma$ is conjugation by $\sigma\in Q$. Hence, $x(C_\sigma-\id)^i\in p\kk$ for all $x\in \kk$ and $i\geq p$. Hence, the generator $T\in\q$ acts on the $\Fp$-module $\Hom(\Lambda^s \kk,\Fp)$, and thus on $\HH^s(\kk,\Fp)$, as $\log \sigma$. In conclusion, the action of $\q$ on $\HH^*(\kk;\F_p)$ is unipotent of class at most $p$ and therefore, we can deduce from Lemma \ref{lem: abelian}, that
	
	\begin{displaymath}
		\HH^r\big(Q;\HH^s(K)\big) \isom \HH^r\big(\q;\HH^s(\kk)\big)
	\end{displaymath}
	as $G/K$-modules. Finally, observe that all the differentials of both spectral sequences are trivial, and hence, $E_{\infty}=E_2$ holds in both cases.
	\end{proof}

\end{document}
